% Generated draft: requires mathematical and source review.
\documentclass[11pt]{article}
\usepackage[margin=25mm]{geometry}
\usepackage{fontspec}
\setmainfont{Latin Modern Roman}
\usepackage{amsmath,amssymb,mathtools,unicode-math}
\setmathfont{Latin Modern Math}
\usepackage{enumitem,needspace,xurl,hyperref}
\setlength{\parindent}{0pt}
\setlength{\parskip}{5pt}
\setlength{\emergencystretch}{3em}
\newcommand{\N}{\mathbb{N}}
\newcommand{\Z}{\mathbb{Z}}
\newcommand{\Q}{\mathbb{Q}}
\newcommand{\R}{\mathbb{R}}
\newcommand{\C}{\mathbb{C}}
\newcommand{\F}{\mathbb{F}}
\newcommand{\A}{\mathbb{A}}
\newcommand{\PP}{\mathbb{P}}
\newcommand{\E}{\mathbb{E}}
\newcommand{\eps}{\varepsilon}
\newcommand{\dd}{\,\mathrm{d}}
\begin{document}
\section*{UnsolvedMath ID: 1865}
\textbf{Source problem code:} catalogue entry
\subsection{Definitions and mathematical statement}
Let $p_1<p_2<\cdots$ be the increasing sequence of prime numbers, and let
\[
 d_n=p_{n+1}-p_n\qquad(n\geq1)
\]
be the consecutive prime gaps. The source literally asks whether there exists $n_0$ such that ``for every $i$ and $n>n_0$'' one has
\[
 d_{n+2i}>d_{n+2i+1}\quad\text{and}\quad
 d_{n+2i+1}<d_{n+2i+2}.
\]
The domain of $i$ is not supplied. If $i$ is interpreted as a nonnegative integer and $n$ independently ranges over every sufficiently large integer, this is not simply a fixed-parity alternation condition: the requirements for adjacent choices of $n$ impose conflicting strict comparisons. The quantifier scope therefore needs explicit correction before it can represent the intended historical question.

\textbf{Natural eventual-alternation interpretation, stated separately.} Is there an integer $N\geq1$ for which
\[
 \forall i\in\Z_{\geq0}:\qquad
 d_{N+2i}>d_{N+2i+1}<d_{N+2i+2}?
\]
This means the entire tail alternates strictly between peaks and troughs, starting at a chosen peak. Allowing the starting index $N$ to vary also covers a tail whose first displayed index would otherwise be a trough. This display fixes a starting index rather than retaining the independently universal $n$ in the import; it is explicitly an interpretation, not an unannounced replacement.

The background attributes to Erdős and Turán the positive lower density of the set $\{n:d_n>d_{n+1}\}$ and records an Erdős offer of \$100.00 for proving that such an eventual-alternation starting index does not exist. These are historical claims supplied by the source, not a statement that a prize is currently available. Its assertion that infinitely many increasing or decreasing triples were unknown is superseded within the record by the dated literature triage.

The triage dated 2026-08-17 reports a theorem of Banks, Freiberg, and Turnage-Butterbaugh producing arbitrarily long strictly increasing runs and arbitrarily long strictly decreasing runs of consecutive prime gaps: for every prescribed length $L$, there is a run
\[
 d_m<d_{m+1}<\cdots<d_{m+L-1},
\]
and likewise a run with all inequalities reversed. The review states that these results refute every natural strict eventual-alternation reading; monotone triples already conflict with alternating strict comparisons. It therefore classifies the question as solved in the literature, in contrast to the registry's imported \emph{open} label, and distinguishes any further quantitative frequency questions from this yes/no question. No proof or independent current literature verification is offered here.
\subsection{Short English statement}
Ask whether consecutive prime gaps eventually alternate strictly up and down. The source's quantifiers are ill-scoped, so a fixed-start interpretation is stated explicitly. Its dated review reports that arbitrarily long increasing and decreasing runs of consecutive prime gaps are known, settling the intended eventual-alternation question negatively.
\subsection{Sources}
\begin{itemize}
\item UnsolvedMath local registry, data/unsolvedmath/problems.json, permanent ID 1865, source code catalogue entry, ``Erdős \$100 Problem on Increasing and Decreasing Gaps,'' including literature triage dated 2026-08-17.
\item Richard Guy, \emph{Unsolved Problems in Number Theory}, Section A11, as identified in the supplied background. No edition, page, or further bibliographic details are supplied; none are inferred. The background attributes earlier density work to Erdős and Turán and the historical \$100.00 offer to Erdős.
\item William D. Banks, Tristan Freiberg, and Caroline L. Turnage-Butterbaugh, ``Consecutive primes in tuples,'' \emph{Acta Arithmetica} 167 (2015), 261--266, \url{https://doi.org/10.4064/aa167-3-4}. The supplied review cites strictly increasing and strictly decreasing runs of every prescribed length. The reference was not independently consulted for this draft.
\item Ulam AI and the UnsolvedMath Contributors, \emph{UnsolvedMath: A Curated Collection of Open Mathematics Problems}. Dataset curation and original metadata are supplied under CC BY 4.0; referenced source material retains its own terms.
\end{itemize}
\end{document}
